\documentclass[10pt,a4paper]{amsart}
\usepackage{amsmath,amssymb,amsthm,graphicx}
% For commutative diagrams we use Xy-pic
\usepackage[all]{xy}

% We use 2cell for 2-commutative diagrams.
\xyoption{2cell}
% Theorem environments.
%
\theoremstyle{plain}
\newtheorem{theorem}[subsection]{Theorem}
\newtheorem{proposition}[subsection]{Proposition}
\newtheorem{lemma}[subsection]{Lemma}

\theoremstyle{definition}
\newtheorem{definition}[subsection]{Definition}
\newtheorem{example}[subsection]{Example}
\newtheorem{exercise}[subsection]{Exercise}
\newtheorem{situation}[subsection]{Situation}

\theoremstyle{remark}
\newtheorem{remark}[subsection]{Remark}
\newtheorem{remarks}[subsection]{Remarks}

\numberwithin{equation}{subsection}

% Macros
%
\def\lim{\mathop{\mathrm{lim}}\nolimits}
\def\colim{\mathop{\mathrm{colim}}\nolimits}
\def\Spec{\mathop{\mathrm{Spec}}}
\def\Hom{\mathop{\mathrm{Hom}}\nolimits}
\def\Ext{\mathop{\mathrm{Ext}}\nolimits}
\def\SheafHom{\mathop{\mathcal{H}\!\mathit{om}}\nolimits}
\def\SheafExt{\mathop{\mathcal{E}\!\mathit{xt}}\nolimits}
\def\Sch{\mathit{Sch}}
\def\Mor{\mathop{\mathrm{Mor}}\nolimits}
\def\Ob{\mathop{\mathrm{Ob}}\nolimits}
\def\Sh{\mathop{\mathit{Sh}}\nolimits}
\def\NL{\mathop{N\!L}\nolimits}
\def\CH{\mathop{\mathrm{CH}}\nolimits}
\def\proetale{{pro\text{-}\acute{e}tale}}
\def\etale{{\acute{e}tale}}
\def\QCoh{\mathit{QCoh}}
\def\Ker{\mathop{\mathrm{Ker}}}
\def\Im{\mathop{\mathrm{Im}}}
\def\Coker{\mathop{\mathrm{Coker}}}
\def\Coim{\mathop{\mathrm{Coim}}}

% Boxtimes
%
\DeclareMathSymbol{\boxtimes}{\mathbin}{AMSa}{"02}

%
% Macros for moduli stacks/spaces
%
\def\QCohstack{\mathcal{QC}\!\mathit{oh}}
\def\Cohstack{\mathcal{C}\!\mathit{oh}}
\def\Spacesstack{\mathcal{S}\!\mathit{paces}}
\def\Quotfunctor{\mathrm{Quot}}
\def\Hilbfunctor{\mathrm{Hilb}}
\def\Curvesstack{\mathcal{C}\!\mathit{urves}}
\def\Polarizedstack{\mathcal{P}\!\mathit{olarized}}
\def\Complexesstack{\mathcal{C}\!\mathit{omplexes}}
% \Pic is the operator that assigns to X its picard group, usage \Pic(X)
% \Picardstack_{X/B} denotes the Picard stack of X over B
% \Picardfunctor_{X/B} denotes the Picard functor of X over B
\def\Pic{\mathop{\mathrm{Pic}}\nolimits}
\def\Picardstack{\mathcal{P}\!\mathit{ic}}
\def\Picardfunctor{\mathrm{Pic}}
\def\Deformationcategory{\mathcal{D}\!\mathit{ef}}
\usepackage{cleveref}
\crefname{originallemma}{original lemma}{original lemmas}
\Crefname{originallemma}{Original lemma}{Original lemmas}
\makeatletter
\newlabel{flat-lemma-char-p}{{0EVW}{source}{}{sourcepointer-1}{}}
\newlabel{flat-lemma-char-p@cref}{{[originallemma][0][]0EVW}{[1][0][]source}{}{}{}}
\makeatother
\crefname{originalreference}{original reference}{original references}
\Crefname{originalreference}{Original reference}{Original references}
\makeatletter
\newlabel{lemma-h-sheaf-colim-F}{{0EVK}{source}{}{sourcepointer-2}{}}
\newlabel{lemma-h-sheaf-colim-F@cref}{{[originalreference][0][]0EVK}{[1][0][]source}{}{}{}}
\makeatother
\crefname{originalsection}{original section}{original sections}
\Crefname{originalsection}{Original section}{Original sections}
\makeatletter
\newlabel{sites-cohomology-section-compare}{{0EWK}{source}{}{sourcepointer-3}{}}
\newlabel{sites-cohomology-section-compare@cref}{{[originalsection][0][]0EWK}{[1][0][]source}{}{}{}}
\makeatother
\title{The h-cohomology of the perfected structure sheaf is the colimit of the cohomologies of the structure sheaf}
\author{The Stacks Project Authors}
\date{stacks-project:etale-cohomology.tex:local-snapshot-20261011; GNU FDL 1.2 or later}
\begin{document}
\maketitle
\noindent\textit{Editorial note.} All delivered mathematical text is the author's own, copied from the pinned Stacks Project chapter file etale-cohomology.tex; no statement or proof step was rewritten, reordered or removed.\par\medskip
\noindent\textit{Editorial note.} The counted proof is the author's complete argument: the proof is written out in full: the identification through the flatness-chapter lemma, the essential-image argument, the Leray chain, and the colimit computation of the fppf cohomology.\par\medskip
\noindent\textit{Editorial note.} The label proposition-h-cohomology-structure-sheaf names the counted statement, and the proof environment immediately following it in the source is delivered verbatim as its proof.\par\medskip
\noindent\textit{Editorial note.} The counted statement is carried by the author's own proposition environment in the chapter source.\par\medskip
\noindent\textit{Editorial note.} The statement is an author Proposition environment, delivered verbatim, including the author's own remark on the notation; the five cited inputs are kept as external pointers with their canonical Stacks tags.\par\medskip
\noindent\textit{Allowed standard prerequisites (editorial).}\begin{itemize}
\item The perfection of the structure sheaf in characteristic p. More on Flatness, the lemma identifying the h-sheafification of the structure sheaf with the perfection colim\_F O. Supplies the identity O\textasciicircum{}h = O\textasciicircum{}perf = colim\_F O.
\item Theperfection lies in the essential image of the h-pushforward. Same chapter, the lemma placing O\textasciicircum{}perf in the essential image of R epsilon\_* from the h-topology. Supplies the object K with R epsilon\_*K = O\textasciicircum{}perf and hence the identification R epsilon\_*O\textasciicircum{}perf = O\textasciicircum{}perf.
\item Leray for the comparison morphism and cohomology of a filtered colimit. Cohomology on Sites, the comparison section for the Leray spectral sequence, together with the computation of fppf cohomology of a filtered colimit. Supplies the chain of identities computing the h-cohomology.
\end{itemize}\medskip
\ifdefined\hypertarget\hypertarget{sourcepointer-1}{}\fi
\noindent\textit{Source reference.} Original lemma 0EVW: sources/stacks-chow/tags.txt, source lines 18601--18601. Prerequisite pointer invoked at proof line(s) 23723; it names a standard fact of a sibling chapter that the counted argument uses, and every step of the counted proof built on it is retained verbatim in the delivered body.\par\medskip
\ifdefined\hypertarget\hypertarget{sourcepointer-2}{}\fi
\noindent\textit{Source reference.} Original reference 0EVK: sources/stacks-etale-cohomology/pinned-original-source.tex, source lines 23651--23651. Prerequisite pointer invoked at proof line(s) 23720; it names a standard fact of the same chapter that the counted argument uses, and every step of the counted proof built on it is retained verbatim in the delivered body.\par\medskip
\ifdefined\hypertarget\hypertarget{sourcepointer-3}{}\fi
\noindent\textit{Source reference.} Original section 0EWK: sources/stacks-chow/tags.txt, source lines 18625--18625. Prerequisite pointer invoked at proof line(s) 23729; it names a standard fact of a sibling chapter that the counted argument uses, and every step of the counted proof built on it is retained verbatim in the delivered body.\par\medskip
\begin{proposition}
\label{proposition-h-cohomology-structure-sheaf}
Let $p$ be a prime number. Let $S$ be a quasi-compact and quasi-separated
scheme over $\mathbf{F}_p$. Then
$$
H^i((\Sch/S)_h, \mathcal{O}^h) =
\colim_F H^i(S, \mathcal{O})
$$
Here on the left hand side by $\mathcal{O}^h$ we mean
the h sheafification of the structure sheaf.
\end{proposition}

\begin{proof}
This is just a reformulation of Lemma \ref{lemma-h-sheaf-colim-F}.
Recall that
$\mathcal{O}^h = \mathcal{O}^{perf} = \colim_F \mathcal{O}$, see
More on Flatness, Lemma \ref{flat-lemma-char-p}.
By Lemma \ref{lemma-h-sheaf-colim-F} we see that
$\mathcal{O}^{perf}$ viewed as an object of $D((\Sch/S)_{fppf})$
is of the form $R\epsilon_*K$ for some $K \in D((\Sch/S)_h)$.
Then $K = \epsilon^{-1}\mathcal{O}^{perf}$ which is actually
equal to $\mathcal{O}^{perf}$ because $\mathcal{O}^{perf}$ is an h sheaf. See
Cohomology on Sites, Section \ref{sites-cohomology-section-compare}.
Hence $R\epsilon_*\mathcal{O}^{perf} = \mathcal{O}^{perf}$
(with apologies for the confusing notation).
Thus the lemma now follows from Leray
$$
R\Gamma_h(S, \mathcal{O}^{perf}) =
R\Gamma_{fppf}(S, R\epsilon_*\mathcal{O}^{perf}) =
R\Gamma_{fppf}(S, \mathcal{O}^{perf})
$$
and the fact that
$$
H^i_{fppf}(S, \mathcal{O}^{perf}) =
H^i_{fppf}(S, \colim_F \mathcal{O}) =
\colim_F H^i_{fppf}(S, \mathcal{O})
$$
as $S$ is quasi-compact and quasi-separated
(see proof of Lemma \ref{lemma-h-sheaf-colim-F}).
\end{proof}

\end{document}
